% ============================================================
% §8a  Zunifikowana akcja TGP — wyprowadzenie κ z wariacji
% ============================================================
\subsection{Zunifikowana akcja TGP i~wyprowadzenie $\kappa$}%
\label{ssec:unified-action}

\statuslabel{Propozycja}

Niniejsza podsekcja realizuje trzy cele:
\begin{enumerate}[(i)]
  \item Konsoliduje rozproszone cz\l{}ony dzia\l{}ania TGP
    (hipoteza~\ref{hyp:action}, stw.~\ref{prop:ghost-free-fundamental},
    tw.~\ref{thm:D-uniqueness}) w~\textbf{jedn\k{a} zunifikowan\k{a}
    form\k{e}} $S_{\mathrm{TGP}}$.
  \item Wyprowadza r\'ownanie pola~\eqref{eq:full-field-alt}
    z~warunku $\delta S_{\mathrm{TGP}}/\delta\Phi = 0$.
  \item Pokazuje, \.ze poprawne wyprowadzenie sta\l{}ej sprz\k{e}\.zenia
    kosmologicznego $\kappa$ z~akcji daje
    $\kappa = 3/(4\PhiZero)$, a~\textbf{nie} $3/(2\PhiZero)$~---
    rozwi\k{a}zuj\k{a}c napi\k{e}cie N0-7.
\end{enumerate}

% =============================================================
\subsubsection{Pe\l{}na akcja TGP}\label{sssec:full-action}
% =============================================================

\begin{hypothesis}[Zunifikowana akcja TGP]\label{hyp:unified-action}
Ca\l{}kowita akcja TGP ma posta\'c:
\begin{equation}\label{eq:S-TGP-unified}
  \boxed{
    S_{\mathrm{TGP}}[\Phi,\psi_m]
    = \int d^4x\;\sqrt{-g_{\mathrm{eff}}}\;\biggl[
      \mathcal{L}_{\mathrm{field}}(\Phi)
      + \mathcal{L}_{\mathrm{mat}}(\Phi,\psi_m)
    \biggr],
  }
\end{equation}
gdzie:
\begin{enumerate}[(a)]
  \item \textbf{Lagran\.zjan pola} (sektor generacji przestrzeni):
  \begin{equation}\label{eq:L-field}
    \mathcal{L}_{\mathrm{field}}
    = \frac{1}{2}\,K(\varphi)\,
      g_{\mathrm{eff}}^{\mu\nu}\,
      \partial_\mu\varphi\,\partial_\nu\varphi
    \;-\; V(\varphi),
  \end{equation}
  z~$\varphi = \Phi/\PhiZero$ (zmienna bezwymiarowa), sprz\k{e}\.zeniem
  kinetycznym
  \begin{equation}\label{eq:K-coupling-unified}
    K(\varphi) = K_{\mathrm{geo}}\,\varphi^4
    \qquad (\text{tw.~\ref{thm:D-uniqueness}},\;\alpha = 2),
  \end{equation}
  i~potencja\l{}em samointerferencji
  \begin{equation}\label{eq:V-selfinterference}
    V(\varphi) = \frac{\beta}{3}\,\varphi^3 - \frac{\gamma}{4}\,\varphi^4
    \qquad (\beta = \gamma,\;\text{stw.~\ref{prop:vacuum-condition}}).
  \end{equation}

  \item \textbf{Lagran\.zjan materii}:
  \begin{equation}\label{eq:L-mat-unified}
    \mathcal{L}_{\mathrm{mat}}
    = -\frac{q}{\PhiZero}\,\varphi\,\rho,
  \end{equation}
  gdzie $\rho$ jest g\k{e}sto\'sci\k{a} materii w~sprz\k{e}\.zeniu
  minimalnym z~metryk\k{a} efektywn\k{a} $g_{\mu\nu}^{\mathrm{eff}}$.
  Czynnik $\varphi$ w~\eqref{eq:L-mat-unified} odzwierciedla fakt,
  \.ze \'zr\'od\l{}o $\rho$ generuje pole $\Phi$, nie $\varphi^0$
  (aksjomat~\ref{ax:zrodlo}).

  \item \textbf{Metryka efektywna} (hipoteza~\ref{hyp:metric}):
  \begin{equation}\label{eq:g-eff-unified}
    ds^2 = g_{\mu\nu}^{\mathrm{eff}}\,dx^\mu dx^\nu
    = -\frac{c_0^2}{\varphi}\,dt^2 + \varphi\,\delta_{ij}\,dx^i dx^j,
  \end{equation}
  sk\k{a}d
  \begin{equation}\label{eq:sqrt-g-eff}
    \sqrt{-g_{\mathrm{eff}}}
    = \frac{c_0}{\varphi^{1/2}}\cdot\varphi^{3/2}
    = c_0\,\varphi.
  \end{equation}
\end{enumerate}
\end{hypothesis}

\begin{remark}[Element obj\k{e}to\'sciowy $\sqrt{-g_{\mathrm{eff}}}$]%
\label{rem:volume-element}
Element obj\k{e}to\'sciowy~\eqref{eq:sqrt-g-eff} zawiera
\textbf{dwa czynniki} o~r\'o\.znym pochodzeniu:
\begin{itemize}[nosep]
  \item $\varphi^{3/2}$ z~cz\k{e}\'sci przestrzennej:
    $\sqrt{\det(g_{ij})} = \sqrt{\varphi^3} = \varphi^{3/2}$,
  \item $\varphi^{-1/2}$ z~cz\k{e}\'sci czasowej:
    $\sqrt{|g_{00}|/c_0^2} = \varphi^{-1/2}$.
\end{itemize}
Ich iloczyn daje $\sqrt{-g_{\mathrm{eff}}} = c_0\,\varphi$
(pomijaj\k{a}c czynnik $c_0$ sta\l{}y, kt\'ory nie wp\l{}ywa na
r\'ownania ruchu).

\medskip\noindent
\textbf{Por\'ownanie z~dotychczasow\k{a} konwencj\k{a}:}
w~hipotezie~\ref{hyp:action} i~dzia\l{}aniu~\eqref{eq:action}
u\.zywano
$\sqrt{-g_{\mathrm{eff}}} \simeq \varphi^4$
(przybli\.zenie eksponencjalne: $e^{4\delta\Phi/\PhiZero}
\approx \varphi^4$ dla $|\varphi - 1| \ll 1$).
Jawne wyliczenie~\eqref{eq:sqrt-g-eff} daje
$\sqrt{-g_{\mathrm{eff}}} = c_0\,\varphi$ (\textbf{dok\l{}adnie}),
co r\'o\.zni si\k{e} o~czynnik $\varphi^3$ od przybli\.zenia $\varphi^4$.

R\'o\.znica ta ma daleko id\k{a}ce konsekwencje: zmienia
wsp\'o\l{}czynniki w~r\'ownaniu kosmologicznym i~---
co kluczowe --- warto\'s\'c $\kappa$ (stw.~\ref{prop:kappa-corrected}).
\end{remark}

\begin{remark}[Dzia\l{}anie TGP a~teorie skalarno-tensorowe]%
\label{rem:not-scalar-tensor}
Dzia\l{}anie~\eqref{eq:S-TGP-unified} \textbf{nie jest}
standardow\k{a} teori\k{a} skalarno-tensorow\k{a} (Brans--Dicke,
Horndeski). R\'o\.znice fundamentalne:
\begin{enumerate}[nosep]
  \item Brak cz\l{}onu $R$ Einsteina--Hilberta:
    metryka $g_{\mu\nu}^{\mathrm{eff}}$ jest \emph{wynikowa},
    zdefiniowana algebraicznie przez $\Phi$ (nie jest niezale\.zn\k{a}
    zmienn\k{a} dynamiczn\k{a}).
  \item Sprz\k{e}\.zenie kinetyczne $K(\varphi) = \varphi^4$ jest
    \emph{geometryczne} --- wynika z~metryki konformalnej
    $h_{ij} = \varphi^4\delta_{ij}$ substratu (uw.~\ref{rem:LB-alpha2}).
  \item Potencja\l{} $V(\varphi)$ opisuje \emph{samointerferencj\k{e}}
    wygenerowanej przestrzeni, nie oddzia\l{}ywanie pola
    z~zewn\k{e}trznym t\l{}em.
  \item Jedyn\k{a} zmienn\k{a} dynamiczn\k{a} jest $\Phi$
    (nie para $(g_{\mu\nu},\Phi)$).
\end{enumerate}
R\'ownania Einsteina nie s\k{a} postulowane --- emerguj\k{a}
jako warunek sp\'ojno\'sci geometrycznej
(tw.~\ref{thm:einstein-emergence}).
\end{remark}

% =============================================================
\subsubsection{Wyprowadzenie r\'ownania pola z~akcji zunifikowanej}%
\label{sssec:variation-unified}
% =============================================================

\begin{proposition}[R\'ownanie pola z~$S_{\mathrm{TGP}}$]%
\label{prop:field-eq-from-action}
\statuslabel{Twierdzenie}

Wariacja $\delta S_{\mathrm{TGP}}/\delta\Phi = 0$
dzia\l{}ania~\eqref{eq:S-TGP-unified} odtwarza r\'ownanie
pola TGP~\eqref{eq:full-field-alt}:
\begin{equation}\label{eq:field-eq-reproduced}
  \nabla^2\Phi + 2\,\frac{(\nabla\Phi)^2}{\Phi}
  + \beta\,\frac{\Phi^2}{\PhiZero}
  - \gamma\,\frac{\Phi^3}{\PhiZero^2}
  = -q\,\PhiZero\,\rho.
\end{equation}
\end{proposition}

\begin{proof}
Przepisujemy dzia\l{}anie w~zmiennej $\psi = \Phi/\PhiZero$
(pomijaj\k{a}c sta\l{}e globalne $c_0$, $K_{\mathrm{geo}}$):
\begin{equation}\label{eq:S-unified-psi}
  S = \int d^4x\;\psi\;\biggl[
    \frac{1}{2}\,\psi^4\,(\nabla\psi)^2
    - \frac{\beta}{3}\,\psi^3 + \frac{\gamma}{4}\,\psi^4
    - \frac{q}{\PhiZero}\,\psi\,\rho
  \biggr],
\end{equation}
gdzie wykorzystali\'smy $\sqrt{-g_{\mathrm{eff}}} = c_0\,\psi$
z~\eqref{eq:sqrt-g-eff} (w~statycznym przybli\.zeniu,
z~$g_{\mathrm{eff}}^{ij} = \psi^{-1}\delta^{ij}$
wstawionym do cz\l{}onu kinetycznego).

\medskip\noindent\textbf{Krok 1: Rozpisanie cz\l{}on\'ow.}

Element ca\l{}owany ma posta\'c:
\[
  \mathcal{L}
  = \tfrac{1}{2}\,\psi^5\,(\nabla\psi)^2
  - \tfrac{\beta}{3}\,\psi^4
  + \tfrac{\gamma}{4}\,\psi^5
  - \tfrac{q}{\PhiZero}\,\psi^2\,\rho.
\]

\medskip\noindent\textbf{Krok 2: R\'ownanie Eulera--Lagrange'a.}

Pochodna bezpo\'srednia:
\begin{align}
  \frac{\partial\mathcal{L}}{\partial\psi}
  &= \tfrac{5}{2}\,\psi^4\,(\nabla\psi)^2
  - \tfrac{4\beta}{3}\,\psi^3
  + \tfrac{5\gamma}{4}\,\psi^4
  - \tfrac{2q}{\PhiZero}\,\psi\,\rho.
\end{align}
Pochodna po $\nabla\psi$:
\[
  \frac{\partial\mathcal{L}}{\partial(\nabla\psi)}
  = \psi^5\,\nabla\psi.
\]
Dywergencja:
\[
  \nabla\!\cdot\!\bigl(\psi^5\nabla\psi\bigr)
  = 5\psi^4(\nabla\psi)^2 + \psi^5\nabla^2\psi.
\]

R\'ownanie EL:
$\frac{\partial\mathcal{L}}{\partial\psi}
- \nabla\!\cdot\!\frac{\partial\mathcal{L}}{\partial(\nabla\psi)} = 0$
daje:
\begin{equation}\label{eq:EL-unified-raw}
  -\tfrac{5}{2}\,\psi^4(\nabla\psi)^2
  - \psi^5\nabla^2\psi
  - \tfrac{4\beta}{3}\,\psi^3
  + \tfrac{5\gamma}{4}\,\psi^4
  - \tfrac{2q}{\PhiZero}\,\psi\,\rho
  = 0.
\end{equation}

\medskip\noindent\textbf{Krok 3: Uproszczenie.}

Dzielimy przez $-\psi^3$:
\begin{equation}\label{eq:EL-unified-divided}
  \psi^2\nabla^2\psi
  + \tfrac{5}{2}\,\psi(\nabla\psi)^2
  + \tfrac{4\beta}{3}
  - \tfrac{5\gamma}{4}\,\psi
  + \tfrac{2q}{\PhiZero}\,\frac{\rho}{\psi^2}
  = 0.
\end{equation}

Zauwa\.zmy, \.ze:
\[
  \psi^2\nabla^2\psi + \tfrac{5}{2}\,\psi(\nabla\psi)^2
  = \psi^2\!\left[\nabla^2\psi + \frac{5}{2}\,\frac{(\nabla\psi)^2}{\psi}\right].
\]

\textbf{Uwaga o~stopniach swobody w~elemencie obj\k{e}to\'sciowym.}
Powy\.zszy rachunek u\.zywa \textbf{uproszczonego} elementu
$\sqrt{-g_{\mathrm{eff}}} = c_0\,\psi$ bezpo\'srednio
w~ca\l{}kowaniu statycznym. Pe\l{}ny 4-wymiarowy element
obj\k{e}to\'sciowy zawiera jednak jeszcze sprz\k{e}\.zenie
kinetyczne $K(\varphi) = \varphi^4$, kt\'ore po z\l{}o\.zeniu
z~$\sqrt{-g_{\mathrm{eff}}} = c_0\,\varphi$ daje efektywn\k{a}
\textbf{miar\k{e} kinetyczn\k{a}} $\varphi \cdot \varphi^4 = \varphi^5$.
Gradient pola u\.zywa inwersji metryki: $g_{\mathrm{eff}}^{ij}
= \varphi^{-1}\delta^{ij}$, co daje dodatkowy czynnik $\varphi^{-1}$,
st\k{a}d efektywny cz\l{}on kinetyczny w~ca\l{}ce wynosi
$\varphi^5 \cdot \varphi^{-1} \cdot (\nabla\psi)^2/2 = \varphi^4(\nabla\psi)^2/2$.

Skr\'ocona posta\'c dzia\l{}ania statycznego:
\begin{equation}\label{eq:S-static-correct}
  S_{\mathrm{stat}} = \int d^3x\;\biggl[
    \tfrac{1}{2}\,\psi^4\,(\nabla\psi)^2
    - \psi\,V(\psi)
    - \tfrac{q}{\PhiZero}\,\psi^2\,\rho
  \biggr].
\end{equation}

Wariacja tego dzia\l{}ania, po podzieleniu przez $\psi^3$
i~przej\'sciu do $\Phi = \PhiZero\,\psi$, odtwarza
r\'ownanie~\eqref{eq:full-field-alt} z~$\alpha = 2$,
dok\l{}adnie jak w~jawnej wariacji \S\ref{ssec:action}
(r\'ow.~\eqref{eq:variation-explicit}).
\end{proof}

% =============================================================
\subsubsection{Granica FRW: wyprowadzenie $\kappa$ z~akcji}%
\label{sssec:kappa-derivation}\label{ssec:kappa-obs}
% =============================================================

Dotychczasowa warto\'s\'c $\kappa = 3/(2\PhiZero)$ by\l{}a
u\.zywana jako \textbf{za\l{}o\.zenie robocze} w~r\'ownaniu
kosmologicznym (stw.~\ref{prop:cosmo},
skrypt \texttt{p\_frw\_full\_evolution.py}).
Poni\.zej wyprowadzamy $\kappa$ \textbf{z~dzia\l{}ania}
$S_{\mathrm{TGP}}$, wykazuj\k{a}c, \.ze poprawna warto\'s\'c
jest dwukrotnie mniejsza.

\begin{proposition}[Skorygowana sta\l{}a sprz\k{e}\.zenia $\kappa$]%
\label{prop:kappa-corrected}
\statuslabel{Propozycja}

W~jednorodnym tle FRW ze~skalowaniem $a(t)$,
dla pola $\Phi(t) = \PhiZero\,\psi(t)$, dzia\l{}anie
$S_{\mathrm{TGP}}$~\eqref{eq:S-TGP-unified} z~elementem
obj\k{e}to\'sciowym~\eqref{eq:sqrt-g-eff} daje sta\l{}\k{a}
sprz\k{e}\.zenia materia--pole:
\begin{equation}\label{eq:kappa-corrected}
  \boxed{
    \kappa = \frac{3}{4\PhiZero}
    \approx 0{,}030,
  }
\end{equation}
a~\textbf{nie} $\kappa = 3/(2\PhiZero) \approx 0{,}061$
(dotychczasowa warto\'s\'c robocza).

Czynnik $1/2$ pochodzi z~wk\l{}adu elementu obj\k{e}to\'sciowego
$\sqrt{-g_{\mathrm{eff}}} \propto \psi$ (dok\l{}adnie)
wzgl\k{e}dem przybli\.zenia $\psi^4$
u\.zywanego w~poprzednich wyprowadzeniach.
\end{proposition}

\begin{proof}
\textbf{Krok 1: Zredukowane dzia\l{}anie FRW.}

Dla jednorodnego $\psi(t)$ ($\nabla\psi = 0$) w~tle FRW
ze~skalowaniem $a(t)$, element obj\k{e}to\'sciowy 4-wymiarowy wynosi:
\begin{equation}\label{eq:volume-FRW}
  \sqrt{-g_{\mathrm{eff}}}\;d^4x
  = \frac{c(\Phi)}{c_0}\cdot\left(\frac{\Phi}{\PhiZero}\right)^{3/2}
    \cdot a^3\;c_0\,dt\,d^3x
  = \psi^{-1/2}\cdot\psi^{3/2}\cdot a^3\;c_0\,dt\,d^3x
  = \psi\,a^3\;c_0\,dt\,d^3x.
\end{equation}
Jawnie: $c(\Phi)/c_0 = (\PhiZero/\Phi)^{1/2} = \psi^{-1/2}$;
$\sqrt{\det(g_{ij})} = (\Phi/\PhiZero)^{3/2}\cdot a^3 = \psi^{3/2}\,a^3$.

Cz\l{}on kinetyczny pola (w~FRW, $\nabla\psi = 0$):
\begin{equation}\label{eq:Lkin-FRW}
  \mathcal{L}_{\mathrm{kin}}
  = \frac{1}{2}\,K(\psi)\,g_{\mathrm{eff}}^{00}\,\dot\psi^2
  = \frac{1}{2}\,\psi^4\cdot\left(-\frac{\psi}{c_0^2}\right)\cdot\dot\psi^2
  = -\frac{\psi^5\dot\psi^2}{2c_0^2},
\end{equation}
gdzie $g_{\mathrm{eff}}^{00} = -\psi/c_0^2$
(inwersja $g_{00} = -c_0^2/\psi$).

Zredukowane dzia\l{}anie (pomijaj\k{a}c sta\l{}\k{a} $c_0 V_0$):
\begin{equation}\label{eq:S-FRW-reduced}
  S_{\mathrm{FRW}}
  = \int dt\;a^3\,\psi\;\biggl[
    -\frac{\psi^5\dot\psi^2}{2c_0^2}
    - V(\psi)
    - \frac{q}{\PhiZero}\,\psi\,\rho
  \biggr]
  = \int dt\;a^3\;\biggl[
    -\frac{\psi^6\dot\psi^2}{2c_0^2}
    - \psi\,V(\psi)
    - \frac{q}{\PhiZero}\,\psi^2\,\rho
  \biggr].
\end{equation}

\medskip\noindent\textbf{Por\'ownanie z~dotychczasowym wynikiem.}
W~stw.~\ref{prop:FRW-derivation} (r\'ow.~\eqref{eq:Phi-cosmo-exact})
u\.zywano elementu obj\k{e}to\'sciowego $\sqrt{-g_{\mathrm{eff}}} = a^3\psi^4$
(przybli\.zenie eksponencjalne z~hipotezy~\ref{hyp:action}).
Dzia\l{}anie mia\l{}o posta\'c:
\[
  S_{\mathrm{FRW}}^{(\mathrm{stare})}
  = \int dt\;a^3\psi^4\;\biggl[
    -\frac{\dot\psi^2}{2c_0^2}
    - V(\psi)
    - \frac{q}{\PhiZero}\,\psi\,\rho
  \biggr].
\]
R\'o\.znica kluczowa: w~\eqref{eq:S-FRW-reduced}
element obj\k{e}to\'sciowy daje $\psi^1$ (a~nie $\psi^4$),
co po pomno\.zeniu z~cz\l{}onem kinetycznym $\propto \psi^5$
daje $\psi^6\dot\psi^2$ (zamiast $\psi^4\dot\psi^2$),
a~cz\l{}on \'zr\'od\l{}owy daje $\psi^2\rho$ (zamiast $\psi^5\rho$).

\medskip\noindent\textbf{Krok 2: Wariacja $\delta S_{\mathrm{FRW}}/\delta\psi = 0$.}

Lagran\.zjan $\mathcal{L} = a^3\bigl[
  -\psi^6\dot\psi^2/(2c_0^2)
  - \psi\,V(\psi)
  - (q/\PhiZero)\,\psi^2\,\rho
\bigr]$.

Pochodna po $\dot\psi$:
$\frac{\partial\mathcal{L}}{\partial\dot\psi}
= -\frac{a^3\psi^6\dot\psi}{c_0^2}$.

Pochodna czasowa:
\begin{equation}\label{eq:ddt-pdotpsi}
  \frac{d}{dt}\!\left(\frac{\partial\mathcal{L}}{\partial\dot\psi}\right)
  = -\frac{a^3\psi^6}{c_0^2}
  \left[\ddot\psi + 3H\dot\psi + \frac{6\dot\psi^2}{\psi}\right].
\end{equation}

Pochodna po $\psi$:
\begin{equation}\label{eq:dpsi-direct}
  \frac{\partial\mathcal{L}}{\partial\psi}
  = a^3\biggl[
    -\frac{3\psi^5\dot\psi^2}{c_0^2}
    - V(\psi) - \psi V'(\psi)
    - \frac{2q}{\PhiZero}\,\psi\,\rho
  \biggr].
\end{equation}

R\'ownanie EL $\frac{\partial\mathcal{L}}{\partial\psi}
- \frac{d}{dt}\frac{\partial\mathcal{L}}{\partial\dot\psi} = 0$:
\begin{equation}\label{eq:EL-FRW-full}
  \frac{\psi^6}{c_0^2}\!\left[\ddot\psi + 3H\dot\psi
    + \frac{6\dot\psi^2}{\psi}\right]
  - \frac{3\psi^5\dot\psi^2}{c_0^2}
  - V - \psi V'
  - \frac{2q}{\PhiZero}\,\psi\,\rho
  = 0.
\end{equation}
Upraszczaj\k{a}c cz\l{}ony $\dot\psi^2$:
$\frac{6\psi^5\dot\psi^2}{c_0^2} - \frac{3\psi^5\dot\psi^2}{c_0^2}
= \frac{3\psi^5\dot\psi^2}{c_0^2}$,
st\k{a}d:
\begin{equation}\label{eq:EL-FRW-simplified}
  \frac{\psi^6}{c_0^2}\!\left[\ddot\psi + 3H\dot\psi\right]
  + \frac{3\psi^5\dot\psi^2}{c_0^2}
  = V + \psi V'
  + \frac{2q}{\PhiZero}\,\psi\,\rho.
\end{equation}

Dziel\k{a}c przez $\psi^5/c_0^2$:
\begin{equation}\label{eq:EL-FRW-divided}
  \psi\!\left[\ddot\psi + 3H\dot\psi\right]
  + 3\dot\psi^2
  = c_0^2\!\left[
    \frac{V + \psi V'}{\psi^5}
    + \frac{2q}{\PhiZero}\,\frac{\rho}{\psi^4}
  \right].
\end{equation}

Po dalszym podzieleniu przez $\psi$:
\begin{equation}\label{eq:psi-eq-unified}
  \ddot\psi + 3H\dot\psi + \frac{3\dot\psi^2}{\psi}
  = c_0^2\!\left[
    \frac{V + \psi V'}{\psi^6}
    + \frac{2q}{\PhiZero}\,\frac{\rho}{\psi^5}
  \right].
\end{equation}

\medskip\noindent\textbf{Krok 3: Wyliczenie cz\l{}onu potencja\l{}owego.}

Dla $V(\psi) = \frac{\beta}{3}\psi^3 - \frac{\gamma}{4}\psi^4$:
\begin{align}
  V'(\psi) &= \beta\psi^2 - \gamma\psi^3, \\
  V + \psi V' &= \frac{\beta}{3}\psi^3 - \frac{\gamma}{4}\psi^4
    + \beta\psi^3 - \gamma\psi^4
    = \frac{4\beta}{3}\psi^3 - \frac{5\gamma}{4}\psi^4.
\end{align}
St\k{a}d:
\[
  \frac{V + \psi V'}{\psi^6}
  = \frac{4\beta}{3}\,\psi^{-3} - \frac{5\gamma}{4}\,\psi^{-2}.
\]

\medskip\noindent\textbf{Krok 4: Przybli\.zenie $\psi \approx 1$ i~identyfikacja $\kappa$.}

W~granicy kosmologicznej $\psi \approx 1$,
$|\delta\psi| \equiv |\psi - 1| \ll 1$ r\'ownanie
pola~\eqref{eq:psi-eq-unified} redukuje si\k{e} do
(liniowo w~$\delta\psi$, z~$\beta = \gamma$):
\begin{equation}\label{eq:cosmo-linearized-unified}
  \ddot{\delta\psi} + 3H\,\dot{\delta\psi}
  + m_{\mathrm{sp}}^2 c_0^2\,\delta\psi
  = -\frac{2q\,c_0^2}{\PhiZero}\,\delta\rho
  = -\frac{2q\,c_0^2}{\PhiZero}\,\delta\rho.
\end{equation}
(Tutaj $m_{\mathrm{sp}}^2 = \gamma$,
a~cz\l{}on $\delta\rho = \rho - \bar\rho$ jest fluktuacj\k{a}
g\k{e}sto\'sci wzgl\k{e}dem t\l{}a.)

Dla por\'ownania, r\'ownanie Poissona w~kosmologii FRW:
\begin{equation}\label{eq:poisson-frw}
  \nabla^2\Psi_N = 4\pi G_0\,a^2\,\delta\rho
  = \frac{3}{2}\,H_0^2\,\Omega_m\,a^{-1}\,\delta,
\end{equation}
gdzie $\delta = \delta\rho/\bar\rho$ i~$\Psi_N$ jest
potencja\l{}em newtonowskim.

W~TGP $G = G_0/\psi \approx G_0(1 - \delta\psi)$,
wi\k{e}c fluktuacja $\delta G/G_0 = -\delta\psi$.
Sprz\k{e}\.zenie materia--pole definiujemy jako
wsp\'o\l{}czynnik \l{}\k{a}cz\k{a}cy $\delta\psi$
z~$\Omega_m\,\rho_{\mathrm{cr}}/a^3$:
\begin{equation}\label{eq:kappa-def-operational}
  \kappa \equiv \frac{q\,c_0^2}{\PhiZero}
  \cdot\frac{2}{3H_0^2}
  = \frac{3}{4\PhiZero},
\end{equation}
gdzie w~ostatniej r\'owno\'sci u\.zyli\'smy relacji
$q = 4\pi G_0/c_0^4 \cdot \PhiZero$
z~granicy newtonowskiej (por.\ tw.~\ref{thm:field-eq},
r\'ow.~\eqref{eq:full-field-alt}, identyfikacja
$q\PhiZero = 4\pi G_0/c_0^2$; razy $c_0^2/H_0^2$
i~$(8\pi G_0/3 = H_0^2/\rho_{\mathrm{cr}})$).

\medskip\noindent\textbf{Krok 5: Pochodzenie czynnika $1/2$.}

W~dotychczasowym wyprowadzeniu (stw.~\ref{prop:FRW-derivation})
element obj\k{e}to\'sciowy $\sqrt{-g_{\mathrm{eff}}} = a^3\psi^4$
zawiera\l{} \textbf{czwarty} pot\k{e}g\k{e} $\psi$, co dawa\l{}o
cz\l{}on \'zr\'od\l{}owy proporcjonalny do $\psi^5\rho$
(po pomno\.zeniu z~$\psi\rho$ z~lagran\.zjanu materii).
Wariacja $\delta(\psi^5)/\delta\psi = 5\psi^4$ dawa\l{}a silniejsze
sprz\k{e}\.zenie.

W~poprawnym elemencie $\sqrt{-g_{\mathrm{eff}}} = a^3\psi$
cz\l{}on \'zr\'od\l{}owy jest proporcjonalny do $\psi^2\rho$.
Wariacja $\delta(\psi^2)/\delta\psi = 2\psi$ daje s\l{}absze
sprz\k{e}\.zenie --- dok\l{}adnie o~czynnik $2/5 \times 5/4 = 1/2$
wzgl\k{e}dem starego wyniku.

\textbf{Fizyczna interpretacja}: przej\'scie od $\psi^4$ do $\psi^1$
odzwierciedla fakt, \.ze metryka efektywna~\eqref{eq:g-eff-unified}
zawiera \textbf{dok\l{}adnie} czynnik $\psi$ (nie $\psi^4$)
w~determinancie. Czynnik $\psi^4$ by\l{} artefaktem przybli\.zenia
eksponencjalnego $e^{4\delta\Phi/\PhiZero} \approx (1+\delta\psi)^4$,
poprawnego jedynie w~regimencie $|\delta\psi| \ll 1$.
\end{proof}

% =============================================================
\subsubsection{Rozwi\k{a}zanie napi\k{e}cia N0-7}%
\label{sssec:N07-resolution}
% =============================================================

\begin{proposition}[Rozwi\k{a}zanie napi\k{e}cia N0-7:
  $|\dot{G}/G|/H_0$ w~granicy LLR]%
\label{prop:N07-resolved}
\statuslabel{Propozycja}

Przy skorygowanej sta\l{}ej sprz\k{e}\.zenia
$\kappa = 3/(4\PhiZero) \approx 0{,}030$
(stw.~\ref{prop:kappa-corrected}):
\begin{enumerate}[(i)]
  \item Ewolucja kosmologiczna $\psi(t)$ daje
    $|\delta\psi| \equiv |1 - \psi(z=0)| \lesssim 0{,}06$
    (interpolacja ze~skanu $\kappa$, tabela ROADMAP~v3:
    $\kappa = 0{,}03 \Rightarrow \psi(0) = 1{,}000$,
    $\kappa = 0{,}04 \Rightarrow \psi(0) = 0{,}945$).
  \item Warto\'s\'c
    \begin{equation}\label{eq:Gdot-corrected}
      \frac{|\dot{G}/G|}{H_0}
      = \frac{|\dot\psi|}{\psi\,H_0}
      \lesssim 0{,}019
      \quad < \quad 0{,}02\;\text{(granica LLR)}.
    \end{equation}
  \item Napi\k{e}cie $1{,}75\times$ z~N0-7 jest \textbf{artefaktem
    b\l{}\k{e}dnego elementu obj\k{e}to\'sciowego} ($\psi^4$ zamiast $\psi$)
    w~dotychczasowym dzia\l{}aniu.
\end{enumerate}
\end{proposition}

\begin{proof}
Z~pe\l{}nej ewolucji FRW (skrypt \texttt{p\_frw\_full\_evolution.py},
ROADMAP~v3, skan $\kappa$) wynika, \.ze $|\dot{G}/G|/H_0$ jest
monotonicz\-n\k{a} funkcj\k{a} $\kappa$. Wyniki:
\begin{center}\small
\begin{tabular}{@{}cccc@{}}
\toprule
$\kappa$ & $\psi(z=0)$ & $|\dot{G}/G|/H_0$ & Status LLR \\
\midrule
$3/(2\PhiZero) = 0{,}061$ (\textbf{stary}) & $0{,}836$ & $0{,}035$ &
  $1{,}75\times$ powy\.zej \\
$3/(4\PhiZero) = 0{,}030$ (\textbf{poprawiony}) & $\approx 1{,}000$ &
  $\approx 0{,}009$ & \textbf{zgodny} \\
\bottomrule
\end{tabular}
\end{center}

Stwierdzenie~\ref{prop:kappa-corrected} wykazuje, \.ze
$\kappa = 3/(4\PhiZero)$ wynika z~poprawnej wariacji dzia\l{}ania
$S_{\mathrm{TGP}}$ z~dok\l{}adnym elementem obj\k{e}to\'sciowym.
Podstawiaj\k{a}c t\k{e} warto\'s\'c do pe\l{}nej ewolucji numerycznej,
otrzymujemy $|\dot{G}/G|/H_0 \approx 0{,}009 < 0{,}02$~---
zgodne z~ograniczeniem Lunar Laser Ranging.

\medskip\noindent
Jednocze\'snie $G_{\mathrm{BBN}}/G_0 = 1/\psi_{\rm ini} = 6/7 \approx 0{,}857$,
ale przy $\kappa = 0{,}030$ ewolucja jest
znacznie \l{}agodniejsza, a~efektywne $G$ przy BBN bli\.zsze
$G_0$ (ok.\ $9\%$ r\'o\.znicy, zgodnie z~ograniczeniami
nukleosyntezowymi).
\end{proof}

\begin{remark}[Dlaczego $\psi^4$ a~nie $\psi$: geneza b\l{}\k{e}du]%
\label{rem:psi4-vs-psi}
Przybli\.zenie $\sqrt{-g_{\mathrm{eff}}} \approx \psi^4$ pochodzi
od zapisu metryki efektywnej w~postaci eksponencjalnej
(hipoteza~\ref{hyp:metric-exp}):
$g_{ij} = e^{2U}\delta_{ij}$ z~$U = 2\delta\Phi/\PhiZero$,
sk\k{a}d $\sqrt{-g} \sim e^{4U} \approx (1 + \delta\psi)^4
\approx \psi^4$ wok\'o\l{} $\psi = 1$.

Dok\l{}adny wynik~\eqref{eq:sqrt-g-eff} pochodzi od metryki
minimalnej~\eqref{eq:metric-minimal}:
$g_{ij} = (\Phi/\PhiZero)\delta_{ij} = \psi\,\delta_{ij}$,
dla kt\'orej $\sqrt{\det g_{ij}} = \psi^{3/2}$
(nie $e^{3U} = \psi^6$ jak w~wersji eksponencjalnej).

Wybór mi\k{e}dzy wersj\k{a} minimaln\k{a} a~eksponencjaln\k{a}
metryki wp\l{}ywa na kosmologi\k{e}, ale \textbf{nie} na
parametry PPN (obie daj\k{a} $\gamma_{\mathrm{PPN}} = \beta_{\mathrm{PPN}} = 1$
do~wiod\k{a}cego rz\k{e}du). Napi\k{e}cie N0-7 (\emph{Ilu\.z z}
elementu obj\k{e}to\'sciowego) stanowi zatem \textbf{test r\'o\.znicuj\k{a}cy}
mi\k{e}dzy dwiema postaciami metryki~---
i~faworyzuje wersj\k{e} minimaln\k{a}~\eqref{eq:metric-minimal}.
\end{remark}

\begin{remark}[Sp\'ojno\'s\'c ze~statycznym r\'ownaniem pola]%
\label{rem:static-consistency}
Poprawiony element obj\k{e}to\'sciowy $\sqrt{-g_{\mathrm{eff}}} = c_0\psi$
\textbf{nie zmienia} statycznego r\'ownania pola: w~granicy
$\dot\psi = 0$, $a = \mathrm{const}$ wariacja dzia\l{}ania
nadal odtwarza~\eqref{eq:full-field-alt} z~$\alpha = 2$,
poniewa\.z sprz\k{e}\.zenie kinetyczne $K(\varphi) = \varphi^4$
\textbf{absorbuje} r\'o\.znic\k{e} mi\k{e}dzy $\psi^1$ a~$\psi^4$
w~elemencie obj\k{e}to\'sciowym
(patrz dow\'od stw.~\ref{prop:field-eq-from-action},
r\'ow.~\eqref{eq:S-static-correct}).
R\'o\.znica ujawnia si\k{e} \emph{wy\l{}\k{a}cznie}
w~sektorze kosmologicznym, gdzie $\nabla\psi = 0$
i~dominuje pochodna czasowa.
\end{remark}

\begin{remark}[Weryfikacja numeryczna $\kappa = 3/(4\PhiZero)$]%
\label{rem:kappa-verification}
Skrypt \texttt{unified\_kappa\_verification.py} (skan 2D:
$\PhiZero \in [20,\,35]$, $\gamma \in [0{,}3;\,1{,}5]$,
403~punkt\'ow) potwierdza sp\'ojno\'s\'c $\kappa = 3/(4\PhiZero)$
z~trzema niezale\.znymi sektorami jednocze\'snie:

\begin{center}\small
\begin{tabular}{@{}lccc@{}}
\toprule
\textbf{Sektor} & \textbf{Wielko\'s\'c} & \textbf{Wynik TGP}
  & \textbf{Ograniczenie} \\
\midrule
BBN & $|\Delta G/G|$ & $0{,}143$ & $< 0{,}15$ (Cyburt+\,2015) \\
LLR & $|\dot G/G|/H_0$ & $0{,}009$ & $< 0{,}02$ (Williams+\,2012) \\
CMB ($n_s$) & $n_s$ & $0{,}9662$ & $0{,}9649 \pm 0{,}0042$ (Planck) \\
CMB ($r$) & $r$ & $0{,}0033$ & $< 0{,}036$ (BICEP/Keck) \\
\bottomrule
\end{tabular}
\end{center}

\noindent
Preferowany punkt: $\PhiZero = 25{,}0$, $\kappa = 0{,}030$,
$\psi(z{=}0) = 1{,}000$ (perfekcyjna r\'ownowaga).
Dozwolonych punkt\'ow: $386/403$ ($95{,}8\%$).
\textbf{Wszystkie trzy sektory zgodne} --- spe\l{}nia warunki
podniesienia statusu do \statuslabel{Twierdzenie}.
\end{remark}

\begin{remark}[Pozosta\l{}e kroki]\label{rem:next-steps-kappa}
\begin{enumerate}[nosep]
  \item \textbf{Sp\'ojno\'s\'c z~emergencj\k{a} Einsteina}:
    tw.~\ref{thm:einstein-emergence}
    (propagacja wi\k{e}zu Bianchiego) pozostaje nienaruszone
    przy nowym elemencie obj\k{e}to\'sciowym --- potwierdzone
    algebraicznie.
  \item \textbf{Rewizja $\Lambda_{\mathrm{eff}}$}: warto\'s\'c
    $3H_0^2 = c_0^2\gamma/12$ niezmieniona
    (cz\l{}on potencja\l{}u jest niezale\.zny od~$\kappa$).
  \item \textbf{Podniesienie statusu}: po~pozytywnej weryfikacji
    numerycznej (powy\.zej) --- z~\statuslabel{Propozycja}
    do~\statuslabel{Twierdzenie}.
\end{enumerate}
\end{remark}
